\section*{Chapter \ref{ch:m1:retail-flow-and-wholesaling} --- Retail Flow and Wholesaling}

\begin{solution}{exo:m1:retail-flow-and-wholesaling:1}
A sale is compared with the best bid: improvement $35.405 - 35.400 = 0.5$
cent a share, \$1.00 on the order. Mid 35.42: \omterm{def:m1:retail-flow-and-wholesaling:effective}{effective half-spread}
$35.420 - 35.405 = 1.5$ cents.
\end{solution}

\begin{solution}{exo:m1:retail-flow-and-wholesaling:2}
For a sale $\epsilon = -1$. Realised: $-(35.405 - 35.412) = +0.7$ cent.
Impact: $-(35.412 - 35.420) = +0.8$ cent. $0.7 + 0.8 = 1.5$: the price moved
the seller's way, against the buyer who provided the liquidity.
\end{solution}

\begin{solution}{exo:m1:retail-flow-and-wholesaling:3}
$0.0018 \times 60\times10^6 \times 252 = \$27.2$ million a year, about \$2.27
per account.
\end{solution}

\begin{solution}{exo:m1:retail-flow-and-wholesaling:4}
$\pi_{\max} = 1.275 - 4p - 0.12$: 0.995 cent at $p = 4\%$, 0.355 cent at
$20\%$, zero at $p = 28.9\%$.
\end{solution}

\begin{solution}{exo:m1:retail-flow-and-wholesaling:5}
B's customer receives 0.30 cent of improvement against 0.10: 0.20 cent a
share better. B's \$1 commission equals that advantage on 500 shares: A is
cheaper for orders below 500 shares, B above. The ``free'' broker is the
expensive one for large orders.
\end{solution}

\begin{solution}{exo:m1:retail-flow-and-wholesaling:6}
Impact 0.05 cent at one second and 0.55 cent at five minutes: the
information in this flow reaches the price slowly --- customers who follow
news or momentum over minutes, not machines. At five minutes the \omterm{def:m1:retail-flow-and-wholesaling:wholesaler}{wholesaler}
keeps 0.20 cent, probably less than its payment and costs: it will give less
improvement on this broker's flow, pay less for it, or decline the most toxic
symbols.
\end{solution}

\begin{solution}{exo:m1:retail-flow-and-wholesaling:7}
The per-trade standard deviation at 300 seconds is about 6.1 cents (the
noise $0.35\sqrt{300}$ dominates), so the standard error from 10\,000 trades
is 0.061 cent. Moving $p$ from 5\% to 8\% lowers the mean by $0.03 \times 3 =
0.09$ cent. Two standard errors below that difference requires $n \ge
(2\times6.1/0.09)^2 \approx 18\,000$ trades: toxicity is measurable per
broker per day, and barely per stock.
\end{solution}

\begin{solution}{exo:m1:retail-flow-and-wholesaling:8}
(i) It is measured against the national best quote, which excludes \omterm{def:m1:us-equity-market-structure:lots}{odd lots}
and hidden or midpoint liquidity and, being taken from the consolidated tape,
is stale when prices move: the true alternative was often better than the
benchmark. (ii) The benchmark spread is itself wider because retail flow has
been removed from exchanges: part of the ``improvement'' returns a cost the
arrangement created. (iii) A gross total scales with volume and spread and
says nothing per share or relative to the realisable alternative, the
midpoint; nor does it net off what the same investors pay indirectly through
the payment to their broker. The relevant statistic is the \omterm{def:m1:retail-flow-and-wholesaling:effective}{effective spread}
as a fraction of the quoted spread, by order size, across \omterm{def:m1:retail-flow-and-wholesaling:wholesaler}{wholesalers}.
\end{solution}

\begin{solution}{pb:m1:retail-flow-and-wholesaling:1}
\textbf{1.} $0.06 \times 3.5 = 0.21$ cent.
\textbf{2.} Effective $1.2 \times 0.8 = 0.96$ cent; realised $0.75$ cent.
\textbf{3.} $0.75 - 0.11 = 0.64$ cent.
\textbf{4.} Gross $0.0075 \times 40\times10^6 = \$300\,000$ a day; margin
$(0.64 - 0.30)/100 \times 40\times10^6 = \$136\,000$.
\textbf{5.} \$75.6 million gross, \$34.3 million of margin.
\textbf{6.} $1.2(1-\theta) = 0.30 + 0.11 + 0.21$ gives $\theta = 48.3\%$.
\textbf{7.} With costs of 0.06: $\theta = 52.5\%$. It outbids you.
\textbf{8.} Your impact falls to 0.13: $\theta = 55.0\%$. You win. Risk
management is a competitive weapon here exactly like cost.
\textbf{9.} Competition is its only lever: the payment being fixed, the
\omterm{def:m1:retail-flow-and-wholesaling:wholesaler}{wholesalers} can compete only on improvement, and the monthly reallocation
turns measured improvement into market share. It also protects the broker's
best-execution defence and avoids dependence on one firm.
\textbf{10.} $0.96 - 1.50 = -0.54$ cent: a loss before costs.
\textbf{11.} Normal stocks: $0.85 \times 40\times10^6 \times 0.0034 =
\$115\,600$. Toxic stocks: margin $-0.54 - 0.11 - 0.30 = -0.95$ cent on 6
million shares, $-\$57\,000$. Total \$58\,600.
\textbf{12.} Effective $2.4$, realised $0.90$ cent; after costs and payment,
$+0.49$ cent. The widening of public quotes restores the \omterm{def:m1:retail-flow-and-wholesaling:wholesaler}{wholesaler}'s
economics: the benchmark does the repricing.
\textbf{13.} Reduce \omterm{def:m1:retail-flow-and-wholesaling:pfof}{price improvement} in those symbols; route more of them to
exchanges instead of internalising; hedge faster and carry less inventory;
ask the broker to renegotiate payment by symbol class. What it cannot do is
refuse individual customers' orders it dislikes while keeping the others.
\textbf{14.} $0.0030 \times 40\times10^6 \times 252 = \$30.2$ million.
\textbf{15.} $0.2 \times 1.2 = 0.24$ cent a share: \$24.2 million.
\textbf{16.} Customers would pay the full half-spread, 1.2 cents; the broker
would pay 0.30 cent a share in taker fees, \$30.2 million a year, instead of
receiving the same amount.
\textbf{17.} 0.96 cent against 1.20: the customers are better off by the
improvement, \emph{assuming} the exchange spread stays at 1.2 cents.
\textbf{18.} With 40 million mostly uninformed shares a day added to the
exchange's flow, $p$ there falls and competition among \omterm{def:m1:what-a-trading-firm-does:market-maker}{market makers} narrows
the quoted spread, possibly by more than the improvement now received. The
comparison cannot be made at today's spread.
\textbf{19.} \textbf{0.64 cent per share.}
\textbf{20.} The right to trade against orders known in advance to carry
little information: a low $p$.
\end{solution}

\begin{solution}{iq:m1:retail-flow-and-wholesaling:1}
A payment from a \omterm{def:m1:what-a-trading-firm-does:market-maker}{market maker} to a broker for routing customer orders to it.
The \omterm{def:m1:what-a-trading-firm-does:market-maker}{market maker} pays because retail orders are, on average, uninformed: the
price moves little against it afterwards, so it keeps most of the spread it
earns, unlike on an exchange where the quoted spread is just enough to cover
\omterm{def:m1:what-a-trading-firm-does:adverse-selection}{adverse selection}. Part of that surplus goes back to the customer as \omterm{def:m1:retail-flow-and-wholesaling:pfof}{price
improvement} and part to the broker.
\iqlookfor{\omterm{def:m1:what-a-trading-firm-does:adverse-selection}{adverse selection} as the reason, and the three-way split.}
\end{solution}

\begin{solution}{iq:m1:retail-flow-and-wholesaling:2}
Effective: signed difference between the execution price and the mid at
execution; it is the customer's cost. Realised: the same against the mid
some time later; it is the liquidity provider's revenue. Their difference is
the price impact, the information in the order.
\iqlookfor{the identity, and ``some time later'' questioned immediately.}
\end{solution}

\begin{solution}{iq:m1:retail-flow-and-wholesaling:3}
Compute the markout curve: share-weighted signed mid change after each fill,
at horizons from milliseconds to tens of minutes, with standard errors, by
symbol class, order size and time of day. Toxic flow has a markout that
grows with the horizon up to the size of the spread or beyond. Measure at the
horizon where the curve flattens --- which also approximates how long the
inventory is actually held before it is hedged. Too short a horizon shows
every flow as benign; too long adds noise without signal.
\iqlookfor{a curve rather than a single number, and standard errors.}
\end{solution}

\begin{solution}{iq:m1:retail-flow-and-wholesaling:4}
The retail order is a one-off with little information; the child order is
one of hundreds from a parent that will keep pushing the price the same way.
On an anonymous book the \omterm{def:m1:what-a-trading-firm-does:market-maker}{market maker} infers the type statistically: order
size patterns, regular timing, persistence of same-side flow, reaction to its
own quotes, venue and order type used. Off-exchange it does not need to
infer: the broker's identity is the label, which is exactly what wholesaling
monetises.
\iqlookfor{autocorrelation of order flow as the signature of a parent
order.}
\end{solution}

\begin{solution}{iq:m1:retail-flow-and-wholesaling:5}
Good: zero commissions and measurable \omterm{def:m1:retail-flow-and-wholesaling:pfof}{price improvement}, fast fills for small
orders, competition among \omterm{def:m1:retail-flow-and-wholesaling:wholesaler}{wholesalers} on execution quality, and a
best-execution duty policed through public reports. Bad: the broker's revenue
depends on a counterparty rather than its customer; improvement is measured
against a benchmark the practice itself widens; segmenting uninformed flow
away from exchanges degrades public price formation and raises costs for
everyone who trades there, including the funds that hold retail savers'
pensions; and order flow concentrates in a few firms.
\iqlookfor{the benchmark problem and the externality on lit markets.}
\end{solution}

\begin{solution}{iq:m1:retail-flow-and-wholesaling:6}
(i) The customers changed: more active or news-driven traders, or a shift
toward volatile symbols --- check impact by symbol class, customer cohort and
time of day. (ii) The broker changed its routing, sending the benign orders
elsewhere and the rest to you --- compare your share of its flow by order
type and the mix against its \omterm{def:m1:retail-flow-and-wholesaling:reports}{Rule 606 report}. (iii) Your own hedging got
worse or slower, so the horizon at which you realise the spread lengthened
--- compare markouts at fixed horizons (unchanged under this hypothesis)
with \omterm{def:m1:pnl-and-positions:realised}{realised P\&L} per share. A fourth: quoted spreads fell while
improvement stayed fixed in cents.
\iqlookfor{hypotheses that make different predictions, each tied to a
dataset.}
\end{solution}
